\section{Experimental Methodology}\label{sec:experiments}
\subsection{Benchmarks and operating grid}
Table~\ref{tab:models} defines the six frozen executable studies. The MLPerf Tiny models use official artifacts and preprocessing~\cite{banbury2021mlperf,mlcommons2026tiny}: CIFAR-10 image classification, Speech Commands keyword spotting~\cite{warden2018speechcommands}, and ToyADMOS anomaly detection~\cite{koizumi2019toyadmos}. DistilBERT/SST-2~\cite{sanh2019distilbert,socher2013sentiment}, Pythia-160M~\cite{biderman2023pythia}, and Qwen3-0.6B~\cite{yang2025qwen3} use validated W8A8 executables, with the language models evaluated on WikiText-2~\cite{merity2017pointer}. W8A8 denotes INT8 weight storage and activation computation; only the declared weight buffers are faultable.

\begin{table}[t]
\caption{Frozen models and evaluation cohorts. Fields are independent fault realizations per BER. Loss units differ by task; language-model perplexity (PPL) refers only to the listed short cohort.}\label{tab:models}
\centering\small
\begin{tabularx}{\linewidth}{l r l r X}\toprule
Model & Eligible bits & Cohort & Fields & Clean loss / quality\\\midrule
ResNet-8 & 618,880 & 1,000 images & 32 & 0.58830 CE / 88.0\% acc.\\
DS-CNN & 176,128 & 1,000 clips & 32 & 0.37939 CE / 90.1\% acc.\\
ToyCar & 2,113,536 & 64 files & 32 & 0.13086 risk / 0.86914 AUC\\
DistilBERT & 344,469,504 & 512 examples & 24 & 0.37473 NLL / 91.41\% acc.\\
Pythia-160M & 988,545,024 & 4,096 tokens & 24 & 3.52956 NLL / 34.109 PPL\\
Qwen3-0.6B & 4,767,875,072 & 4,096 tokens & 24 & 2.82863 NLL / 16.922 PPL\\\bottomrule
\end{tabularx}
\Description{Six models, eligible storage populations, fixed input cohorts, fields per BER and clean metrics.}
\end{table}

All studies evaluate 6.25\%, 12.5\%, 18.75\%, and 25\% protection plus zero/full controls. BERs are $10^{-5},10^{-4},3\!\times\!10^{-4},10^{-3},3\!\times\!10^{-3},10^{-2},3\!\times\!10^{-2},10^{-1}$; ToyCar uses the subset $10^{-5},10^{-4},10^{-3},3\!\times\!10^{-3},10^{-2}$. Every grid includes the six structural/gradient methods in Table~\ref{tab:methods}; FKeras is measured on the three TFLite models only. No matched PrioriFI-derived allocation is available.

Ranking uses 64 calibration inputs for ResNet-8, KWS and DistilBERT, 16 normal training files for ToyCar, and 8,184 validation targets for each language model. Final inputs are separate within each study, but several cohorts have prior exploratory exposure; they are not untouched confirmation samples. Each model has one checkpoint. TFLite random allocation uses four mask seeds per budget; other methods, and all transformer methods, use one frozen mask per budget.

\subsection{Metrics and uncertainty}
Classification uses accuracy and cross-entropy, in natural-log units (nats). TFLite probabilities are floored at $10^{-12}$. Transformer NLL is capped at 30 nats per example/token, also the penalty for a failed or nonfinite unit. ToyCar evaluates every window under the same fault state, aggregates one score per complete file, and computes ROC-AUC with half credit for ties. Its task risk is $1-\mathrm{AUC}$, not reconstruction MSE.

Contrasts are paired by proposed field and retain their original Holm comparison families. Confidence intervals quantify fault randomness conditional on the fixed model, inputs and masks; they do not replicate training or input sampling. Transformer contrasts use paired bootstrap intervals. The additional Qwen crossover study uses 24 new fields at each of 5\% and 7.5\% BER; its intervals are exploratory.

Pairing asks how much two allocations differ under the same proposed corruption, rather than comparing two unrelated trial averages. A wide distribution of fault outcomes can coexist with a precise paired difference if the methods respond similarly to easy and difficult fields. Conversely, overlapping mean intervals alone establish neither equivalence nor changing method preference. We therefore use the paired contrasts to support reversals and the full curves to show their operating context.

The isolated-damage study uses 320 common stratified candidate bits per TFLite model; its correlations are sample-limited. The conditional KWS test uses 32 independently sampled six-bit pools within a 12.5\% absolute-gradient mask, preserving two protected bits per pool. At BER 5\%, each pool uses eight selection fields and 32 independent final master fields on 384 clips. Its primary contrast is conditional versus exact-isolated selection; two secondary contrasts use Holm correction. Pool-level uncertainty is distinct from field-level uncertainty.

All reported measurements come from finalized, checksum-linked artifacts. Model/cohort identities, masks, schedules, detailed numerical contracts and full grids accompany the paper.
